\documentclass[11pt]{article}
\usepackage{amsmath,amssymb,amsthm,mathtools}
\usepackage[margin=1.05in]{geometry}
\usepackage{hyperref}
\usepackage{booktabs}

\newtheorem{theorem}{Theorem}[section]
\newtheorem{proposition}[theorem]{Proposition}
\newtheorem{lemma}[theorem]{Lemma}
\newtheorem{corollary}[theorem]{Corollary}
\theoremstyle{definition}
\newtheorem{remark}[theorem]{Remark}
\newtheorem{example}[theorem]{Example}

\newcommand{\hgt}{\operatorname{ht}}
\newcommand{\Q}{\mathbb{Q}}
\newcommand{\Z}{\mathbb{Z}}
\newcommand{\Lam}{\Lambda}
\newcommand{\wtB}{\widetilde{B}}

\title{Two combinatorial gaps closed:\\
connectivity and odd cycles in the reflection graph $G(\mu)$}
\author{Clio}
\date{11 September 2026}

\begin{document}
\maketitle

\begin{abstract}
Let $G(\mu)$ be the labelled reflection graph attached to a partition $\mu\vdash n$ in
\cite[\S6]{Q140}, whose vertices are the $n$ columns of the local-identity matrix
$M^{(\mu)}(t)$ and whose edges are the two-term relations of that paper's
Proposition~6.2. Two statements about $G(\mu)$ were left open there, and together they
finish that paper's classification theorem. \textbf{Both are proved here, and in a
sharper form than was asked for.}
\begin{itemize}
\item \textbf{(G1)} $G(\mu)$ is connected for every $\mu$ --- indeed it has diameter
      at most $4$, with the corner cell $(1,1)$ at distance $\le2$ from every vertex.
\item \textbf{(G2)} For $n\ge6$, $G(\mu)$ carries a cycle of nontrivial label product.
      What is proved is stronger and exponent-free: $G(\mu)$ contains a
      \emph{triangle} --- hence an odd cycle, whose label product is automatically
      $\ne1$ --- for \emph{every} $\mu\vdash n$ with $n\ge4$ except exactly
      $\mu\in\{(2,2),(3,2),(2,2,1)\}$. Triangle-freeness is classified completely
      (Theorem~\ref{thm:tri}).
\end{itemize}
The engine is a reformulation (Theorem~\ref{thm:struct}): $G(\mu)$ lives on the Young
diagram of $\mu$, its edges run only along rows and columns, and
\[
   (i,k)\sim(i,m)\iff h_{ik}+h_{im}\notin\{\text{hook lengths of row }i\},\qquad
   (i,k)\sim(i',k)\iff h_{ik}+h_{i'k}\notin\{\text{hook lengths of column }k\}.
\]
The variable $t$ disappears from the combinatorics entirely. We also supply the one
missing proof upstream: \cite[Prop.~6.2]{Q140}, the reflection relations themselves, were
stated there with a sketch and verified to $n\le8$; \S\ref{sec:reflect} proves them, by
exhibiting the witnessing row of $M^{(\mu)}$ explicitly and showing it has exactly two
nonzero entries. Consequently \cite[Thm.~3.1]{Q140} --- the classification of the $\mu$ for
which the Khanna--Loehr local identity with $T=S$ is consistent over $\Q(t)$ --- is now
\textbf{proved for all $\mu$ and all $n$}, not verified to $n\le9$, and no step of the chain
rests on a computation.

We also \textbf{correct} \cite[\S7(2)]{Q140}: the obstruction recorded there for
$\mu=(3,1,1,1,1)$ (``no second edge inside row~$1$ exists'') is false. It arises from
using a sufficient condition for a type-(II) edge as if it were necessary. Row~$1$ of
$(3,1,1,1,1)$ does carry a triangle, exhibited below and confirmed independently from
the matrix.
\end{abstract}

\section{Setting and notation}

We use the notation of \cite{Q140} and recall only what is needed. Fix $\mu\vdash n$ with
$\ell=\ell(\mu)$ parts and conjugate $\mu'$. Work in the $\Z$-Maya convention: the bead set is
\[
   B=B(\mu)=\{\,b_j:=\mu_j-j\ :\ j\ge1\,\},\qquad \mu_j:=0 \text{ for } j>\ell,
\]
so $b_1>b_2>\cdots$ strictly, and $b_j=-j$ for $j>\ell$. Let $U=\Z\setminus B$ be the set of
\emph{holes}. The cell $(i,j)$ of $\mu$ has hook length
$h_{ij}=\mu_i-j+\mu'_j-i+1$.

By \cite[Prop.~2.1]{Q140} the columns of $M^{(\mu)}(t)$ --- equivalently the elements of
$C(\mu)=\{\gamma:\mu/\gamma \text{ a rim hook}\}$ --- are in bijection with the pairs
$(b,b')$ with $b\in B$, $b'\in U$, $b'<b$, the pair $(b,b')$ corresponding to moving the bead
$b$ down to the hole $b'$; and $|C(\mu)|=n$. Write $w(b,b')$ for the unknown attached to that
column. The edges of $G(\mu)$ are the two-term relations of \cite[Prop.~6.2]{Q140}, which we
quote verbatim:

\begin{proposition}[{\cite[Prop.~6.2]{Q140}}]\label{prop:reflect}
\begin{enumerate}
\item[\textup{(I)}] Let $c<b$ be beads and $b'<c$ a hole, and suppose $e:=b+c-b'$ is a hole
(it is automatically $>b$). Then $M$ has a row with exactly the two entries in columns
$(b,b')$ and $(c,b')$, giving $w(c,b')=-t^{\,|B\cap(c,b)|}w(b,b')$.
\item[\textup{(II)}] Let $b$ be a bead and $b'<e<b$ holes, and suppose $c:=b'+e-b$ is a bead
(it is automatically $<b'$). Then $M$ has a row with exactly the two entries in columns
$(b,b')$ and $(b,e)$, giving $w(b,e)=-t^{\,1+|B\cap(b',e)|}w(b,b')$.
\end{enumerate}
\end{proposition}

Each edge is labelled by its ratio $-t^{h}\in\Q(t)^\times$. The criterion to be fed is:

\begin{theorem}[{\cite[Thm.~6.3]{Q140}}]\label{thm:graph}
If every connected component of $G(\mu)$ contains a cycle whose product of edge labels
(traversed with signs) is $\ne1$ in $\Q(t)$, then $w=0$ is forced by the rows $\lambda\ne\mu$,
and the local identity \eqref{eq:local} is inconsistent.
\end{theorem}

\begin{equation}\label{eq:local}
  \sum_{L=1}^{n}\ \sum_{\gamma\in S(\lambda,L)\cap S(\mu,L)}
     t^{\hgt(\lambda/\gamma)}\,\wtB(\mu,\gamma)\;=\;\delta_{\lambda\mu},
  \qquad \forall\,\lambda\vdash n .
\end{equation}

\section{Proposition~\ref{prop:reflect} itself, proved}\label{sec:reflect}

Proposition~\ref{prop:reflect} is stated in \cite{Q140} with a sketch (``sorting the four
positions shows that such a $\lambda$ \emph{generically} admits exactly two decompositions'')
and is verified there for $n\le8$; it is not proved. Since it is the only link between
$G(\mu)$ and the matrix $M^{(\mu)}$, we prove it here, so that the chain ending in
Corollary~\ref{cor:class} contains no computed step. We use one standard fact, proved in
\cite[Prop.~2.1]{Q140}: if $B(\lambda)=B(\gamma)\setminus\{d\}\cup\{d+L\}$ with $d\in B(\gamma)$
and $d+L\notin B(\gamma)$, then $\lambda/\gamma$ is a rim hook of size $L$ and
$\hgt(\lambda/\gamma)=|B(\gamma)\cap(d,d+L)|$; and conversely every rim hook arises so.
(The open interval contains neither endpoint, so $B(\gamma)$ may be replaced by $B(\lambda)$.)

\begin{proof}[Proof of Proposition~\ref{prop:reflect}]
Both cases are the same configuration, and we treat them together. Let $\{b,c\}\subset B$ be
two beads and $\{b',e\}\subset U$ two holes with
\[
   b+c=b'+e .
\]
Reflection in the common midpoint exchanges $b\leftrightarrow c$ and $b'\leftrightarrow e$, so
after naming $c<b$ and $b'<e$ exactly two orderings are possible: either $b'<c<b<e$ (the beads
are nested inside the holes) --- this is case (I) --- or $c<b'<e<b$ (the holes nested inside
the beads) --- case (II). Interleavings are impossible: $b'<c$ and $e<b$ would give
$b'+e<b+c$, and $c<b'$ with $b<e$ would give $b+c<b'+e$.

Define
\begin{equation}\label{eq:lam}
   B_\lambda \;:=\; \bigl(B\setminus\{b,c\}\bigr)\cup\{b',e\},
\end{equation}
a bead set with the same cofinite left tail as $B$; let $\lambda$ be the corresponding
partition. Since $b+c=b'+e$, the sum of the positions is unchanged, so $|\lambda|=|\mu|=n$;
and $\lambda\ne\mu$ because $b\in B\setminus B_\lambda$.

\smallskip\noindent\emph{Step 1: the two columns.}
In case (I) put $\gamma_1=(b,b')$ and $\gamma_2=(c,b')$; in case (II) put $\gamma_1=(b,b')$
and $\gamma_2=(b,e)$. In both cases $\gamma_1,\gamma_2\in C(\mu)$ by definition (a bead moved
down to a hole below it: $b'<b$ always, $b'<c$ in case (I), $e<b$ in case (II)). Writing
$B_1=B\setminus\{b\}\cup\{b'\}$ we get, in either case,
\[
   B_\lambda \;=\; B_1\setminus\{c\}\cup\{e\},
\]
an up-move of the bead $c$ by $e-c=b-b'=|\mu|-|\gamma_1|$. So $\lambda/\gamma_1$ is a rim hook
of the right size and $M_{\lambda\gamma_1}\ne0$. For $\gamma_2$:
in case (I), $B_2=B\setminus\{c\}\cup\{b'\}$ and $B_\lambda=B_2\setminus\{b\}\cup\{e\}$, an
up-move of $b$ by $e-b=c-b'=|\mu|-|\gamma_2|$;
in case (II), $B_2=B\setminus\{b\}\cup\{e\}$ and $B_\lambda=B_2\setminus\{c\}\cup\{b'\}$, an
up-move of $c$ by $b'-c=b-e=|\mu|-|\gamma_2|$ (note $c<b'$ here). Either way
$M_{\lambda\gamma_2}\ne0$.

\smallskip\noindent\emph{Step 2: there are no other nonzero entries.}
Let $\gamma\in C(\mu)$ with $M_{\lambda\gamma}\ne0$, say $B(\gamma)=B\setminus\{x\}\cup\{y\}$
with $x\in B$, $y\in U$, $y<x$. Then $\lambda/\gamma$ is a rim hook, so $B_\lambda$ and
$B(\gamma)$ differ by a single bead move; in particular
$|B(\gamma)\setminus B_\lambda|=1$. Now an element of $B(\gamma)$ lying outside $B_\lambda$ is
either an element of $B\setminus\{x\}$ lying in $\{b,c\}$, or $y$ itself when
$y\notin\{b',e\}$ (as $y\notin B$). Hence
\[
   B(\gamma)\setminus B_\lambda \;=\; \bigl(\{b,c\}\setminus\{x\}\bigr)\ \cup\
   \bigl(\{y\}\ \text{if } y\notin\{b',e\}\bigr).
\]
If $x\notin\{b,c\}$ this set already has two elements, a contradiction. So $x\in\{b,c\}$, and
then $\{b,c\}\setminus\{x\}$ is a singleton, forcing $y\in\{b',e\}$. Four pairs $(x,y)$ remain,
and the constraint $y<x$ eliminates two of them:
\[
\begin{array}{llll}
\text{case (I)}\ \ (b'<c<b<e): & (b,b')\ \checkmark & (c,b')\ \checkmark &
   (b,e),(c,e)\ \text{fail since } e>b>c;\\[2pt]
\text{case (II)}\ (c<b'<e<b): & (b,b')\ \checkmark & (b,e)\ \checkmark &
   (c,b'),(c,e)\ \text{fail since } c<b'<e.
\end{array}
\]
So $\gamma\in\{\gamma_1,\gamma_2\}$: the row $\lambda$ of $M$ has exactly two nonzero entries.

\smallskip\noindent\emph{Step 3: the exponents.}
Since $\lambda\ne\mu$, the $\lambda$-equation of \eqref{eq:local} reads
$t^{\hgt_1}w(\gamma_1)+t^{\hgt_2}w(\gamma_2)=0$, i.e.\
$w(\gamma_2)=-t^{\,\hgt_1-\hgt_2}w(\gamma_1)$, where $\hgt_j=\hgt(\lambda/\gamma_j)$.

\emph{Case (I).} $\hgt_1=|B_1\cap(c,e)|$ with $B_1=B\setminus\{b\}\cup\{b'\}$; the open
interval $(c,e)$ contains $b$ but not $b'$ (as $b'<c$), so $\hgt_1=|B\cap(c,e)|-1$. And
$\hgt_2=|B_2\cap(b,e)|$ with $B_2=B\setminus\{c\}\cup\{b'\}$; the interval $(b,e)$ contains
neither $c$ nor $b'$, so $\hgt_2=|B\cap(b,e)|$. Hence, using $b\in B$,
\[
   \hgt_1-\hgt_2=|B\cap(c,e)|-1-|B\cap(b,e)|=|B\cap(c,b]|-1=|B\cap(c,b)| ,
\]
giving $w(c,b')=-t^{\,|B\cap(c,b)|}w(b,b')$.

\emph{Case (II).} $\hgt_1=|B_1\cap(c,e)|$; here $c<b'<e<b$, so $(c,e)$ contains $b'$ but not
$b$, whence $\hgt_1=|B\cap(c,e)|+1$. And $\hgt_2=|B_2\cap(c,b')|$ with
$B_2=B\setminus\{b\}\cup\{e\}$; the interval $(c,b')$ contains neither $b$ nor $e$, whence
$\hgt_2=|B\cap(c,b')|$. Hence, using that $b'$ is a hole,
\[
   \hgt_1-\hgt_2=|B\cap(c,e)|+1-|B\cap(c,b')|=1+|B\cap[b',e)|=1+|B\cap(b',e)| ,
\]
giving $w(b,e)=-t^{\,1+|B\cap(b',e)|}w(b,b')$. \qedhere
\end{proof}

\begin{remark}
The proof gives the witnessing row explicitly: \eqref{eq:lam}. It does \emph{not} prove the
completeness statement of \cite[\S6]{Q140} --- that (I) and (II) exhaust the two-term rows of
$M$ --- and completeness is not needed for Theorem~\ref{thm:graph}, which only requires each
edge of $G(\mu)$ to be a genuine relation.
\end{remark}

\section{The structure theorem: $G(\mu)$ is a hook-length graph}\label{sec:struct}

Everything below rests on the following dictionary. It is implicit in the proof of
\cite[Prop.~2.1]{Q140} but is not stated there, and stating it is what makes the two gaps
easy.

\begin{lemma}\label{lem:dict}
Let $u_1<u_2<\cdots$ enumerate the holes of $B(\mu)$ in increasing order. Then:
\begin{enumerate}
\item $u_1=-\ell$; every position $<-\ell$ is a bead, and every position $>b_1$ is a hole.
\item For each $i\le\ell$, the holes below $b_i$ are exactly $u_1<\cdots<u_{\mu_i}$, and
      \[ h_{ik}=b_i-u_k \qquad (1\le k\le \mu_i). \]
\item Consequently $u_k<b_i \iff k\le\mu_i \iff i\le\mu'_k$, and the map
      $(i,k)\mapsto(b_i,u_k)$ is a bijection from the cells of $\mu$ to the vertices of
      $G(\mu)$.
\end{enumerate}
\end{lemma}

\begin{proof}
(1) The beads are $b_j=\mu_j-j$. For $j>\ell$ we have $b_j=-j\le-\ell-1$, so every position
$\le-\ell-1$ is a bead. The position $-\ell$ is a hole: $b_j=-\ell$ would need $\mu_j=j-\ell$,
which is $\le0$ for $j\le\ell$ and negative for $j>\ell$. Since $b_1$ is the largest bead,
every position $>b_1$ is a hole.

(2) This is the classical first-column hook-length identity, proved as part of
\cite[Prop.~2.1]{Q140}: the set $\{v<b_i : v\notin B\}$ equals $\{b_i-h_{ij}:1\le j\le\mu_i\}$.
As $h_{i1}>h_{i2}>\cdots>h_{i\mu_i}$ strictly, the values $b_i-h_{ij}$ increase strictly with
$j$, so listing them in increasing order matches the global enumeration: $b_i-h_{ij}=u_j$.
(The count also confirms $u_1=b_i-h_{i1}=-\ell$ for every $i$, since
$h_{i1}=\mu_i+\ell-i$ and $b_i=\mu_i-i$.)

(3) Immediate from (2): there are exactly $\mu_i$ holes below $b_i$, namely $u_1,\dots,u_{\mu_i}$;
and $k\le\mu_i\iff i\le\mu'_k$ is the definition of the conjugate. Summing over $i$ gives $n$
pairs, matching $|C(\mu)|=n$.
\end{proof}

We henceforth name vertices by cells. Write
\[
   H_i=\{h_{i1},\dots,h_{i\mu_i}\} \quad(\text{hooks of row } i), \qquad
   H^{(k)}=\{h_{1k},\dots,h_{\mu'_k,k}\}\quad(\text{hooks of column } k).
\]

\begin{theorem}[structure]\label{thm:struct}
The vertex set of $G(\mu)$ is the set of cells of $\mu$, and every edge joins two cells in the
same row or two cells in the same column. Precisely, for $k\ne m$ and $i\ne i'$:
\begin{align}
  (i,k)\sim(i,m) \ \text{\textup{(type II)}} &\iff h_{ik}+h_{im}\notin H_i, \label{eq:row}\\
  (i,k)\sim(i',k)\ \text{\textup{(type I)}} &\iff h_{ik}+h_{i'k}\notin H^{(k)}. \label{eq:col}
\end{align}
There are no other edges. The label of the edge \eqref{eq:row} is $-t^{\,1+|B\cap(u_k,u_m)|}$
and of \eqref{eq:col} is $-t^{\,|i'-i|-1}$.
\end{theorem}

\begin{proof}
A type-(II) relation fixes the bead $b$ and varies the hole; by Lemma~\ref{lem:dict} that is
two cells in the same row. A type-(I) relation fixes the hole $b'$ and varies the bead: two
cells in the same column. As (I) and (II) are the only relations in
Proposition~\ref{prop:reflect}, there are no other edges.

\emph{Type \textup{(II)}.} Take $b=b_i$ and holes $u_k<u_m<b_i$, i.e.\ $1\le k<m\le\mu_i$.
Using $u_k=b_i-h_{ik}$ and $u_m=b_i-h_{im}$,
\[
   c \;=\; u_k+u_m-b_i \;=\; b_i-(h_{ik}+h_{im}).
\]
Since $h_{im}\ge1$ we have $c<u_k<b_i$, so $c$ is a hole if and only if $c$ is one of the holes
below $b_i$, i.e.\ $c=b_i-h_{iq}$ for some $q\le\mu_i$, i.e.\ $h_{ik}+h_{im}\in H_i$. (If
$h_{ik}+h_{im}>h_{i1}$ then $c<-\ell$ and $c$ is a bead by Lemma~\ref{lem:dict}(1); this is the
case $h_{ik}+h_{im}\notin H_i$ as well.) The edge exists exactly when $c\in B$, giving
\eqref{eq:row}.

\emph{Type \textup{(I)}.} Take beads $b=b_i>c=b_{i'}$, so $i<i'$, and a hole $u_k<b_{i'}$,
i.e.\ $k\le\mu_{i'}$, equivalently $i'\le\mu'_k$. Then
\[
   e \;=\; b_i+b_{i'}-u_k \;=\; u_k+h_{ik}+h_{i'k},
\]
and $e>b_i$ since $h_{i'k}\ge1$. Every bead is some $b_j$; and $e>u_k\ge-\ell$ forces $j\le\ell$,
while $e>u_k$ forces $j\le\mu'_k$ by Lemma~\ref{lem:dict}(3). So $e$ is a bead if and only if
$u_k+h_{ik}+h_{i'k}=u_k+h_{jk}$ for some $j\le\mu'_k$, i.e.\ $h_{ik}+h_{i'k}\in H^{(k)}$. The
edge exists exactly when $e\in U$, giving \eqref{eq:col}. For its label,
$|B\cap(b_{i'},b_i)|=|\{b_{i+1},\dots,b_{i'-1}\}|=i'-i-1$.
\end{proof}

\begin{remark}
Theorem~\ref{thm:struct} is manifestly self-conjugate: conjugation of $\mu$ acts on the abacus
by $x\mapsto -1-x$, exchanging beads with holes, hence exchanging type (I) with type (II), and
on cells by transposition, which exchanges \eqref{eq:row} with \eqref{eq:col}. Every statement
below therefore comes in a row/column pair, and we prove only one of each.
\end{remark}

\section{(G1): connectivity}

\begin{lemma}[hubs]\label{lem:hub}
For every $i$ and every $2\le m\le\mu_i$ the edge $(i,1)\sim(i,m)$ is present; for every $k$
and every $2\le i'\le\mu'_k$ the edge $(1,k)\sim(i',k)$ is present.
\end{lemma}

\begin{proof}
Hook lengths decrease strictly along a row, so $h_{i1}=\max H_i$. Hence
$h_{i1}+h_{im}>h_{i1}=\max H_i$, so $h_{i1}+h_{im}\notin H_i$, and \eqref{eq:row} gives the
edge. Hook lengths decrease strictly down a column, so $h_{1k}=\max H^{(k)}$ and the same
one-line argument with \eqref{eq:col} gives the second statement.
\end{proof}

\begin{theorem}[(G1)]\label{thm:G1}
$G(\mu)$ is connected for every partition $\mu$. More precisely every vertex is within
distance $2$ of the corner cell $(1,1)$, so $\operatorname{diam} G(\mu)\le4$.
\end{theorem}

\begin{proof}
Let $(i,k)$ be a cell. By Lemma~\ref{lem:hub} applied to row $i$, $(i,k)\sim(i,1)$ (or
$k=1$). By Lemma~\ref{lem:hub} applied to column $1$, $(i,1)\sim(1,1)$ (or $i=1$). So
$d\bigl((i,k),(1,1)\bigr)\le2$.
\end{proof}

\noindent This is exactly the ``hub'' idea sketched in \cite[\S7(2)]{Q140}, but with the
\emph{exact} edge condition \eqref{eq:row} in place of the sufficient condition used there,
and applied to every row rather than only row~$1$.

\section{Odd cycles have nontrivial label product}

\begin{lemma}\label{lem:odd}
Every edge label of $G(\mu)$ has the form $-t^{h}$ with $h\in\Z$. If a cycle has odd length
$\ell$, its label product is $-t^{m}$ for some $m\in\Z$, hence $\ne1$ in $\Q(t)$.
\end{lemma}

\begin{proof}
In $\Q(t)^\times$ we have $(-t^{h})(-t^{-h})=t^{0}=1$, so $(-t^{h})^{-1}=-t^{-h}$: traversing
an edge backwards inverts the exponent but \emph{does not} change the sign. Hence each of the
$\ell$ edges of the cycle contributes exactly one factor $-1$, and the product is
$(-1)^{\ell}t^{\sum_j\pm h_j}$. For $\ell$ odd this is $-t^{m}$. If $-t^{m}=1$ in $\Q(t)$,
evaluate at $t=1$ to get $-1=1$, a contradiction.
\end{proof}

\begin{corollary}\label{cor:nonbip}
If $G(\mu)$ is connected and not bipartite, the hypothesis of Theorem~\ref{thm:graph} holds and
\eqref{eq:local} is inconsistent over $\Q(t)$. \emph{No exponent bookkeeping is required.}
\end{corollary}

\begin{proof}
A connected non-bipartite graph contains an odd cycle; apply Lemma~\ref{lem:odd} and
Theorem~\ref{thm:G1}.
\end{proof}

\section{(G2): classification of triangle-free $G(\mu)$}

\begin{lemma}[line lemma]\label{lem:line}
Let $h_1>h_2>\cdots>h_r\ge1$ be integers and put $S=\{h_1,\dots,h_r\}$. Call a pair
$2\le q<s\le r$ \emph{good} if $h_q+h_s\notin S$. Then:
\begin{enumerate}
\item If $\{q,s\}$ is good, the three indices $\{1,q,s\}$ are pairwise joined in the sense of
      \eqref{eq:row} \textup{(}resp.\ \eqref{eq:col}\textup{)}.
\item A good pair exists if and only if $r\ge4$, or $r=3$ and $h_2+h_3\ne h_1$.
\end{enumerate}
\end{lemma}

\begin{proof}
(1) The pairs $\{1,q\}$ and $\{1,s\}$ are good by the argument of Lemma~\ref{lem:hub}
($h_1+h_x>h_1=\max S$), and $\{q,s\}$ is good by hypothesis.

(2) First, if $h_q+h_s\in S$, say $h_q+h_s=h_j$, then $h_j>h_q$ forces $j<q$; in particular for
$q=2$ we get $j=1$, i.e.\ $h_2+h_s=h_1$.

If $r\le2$ there is no pair at all. If $r=3$ the only pair is $\{2,3\}$, good iff
$h_2+h_3\ne h_1$. If $r\ge4$, consider the two pairs $\{2,3\}$ and $\{2,4\}$. Were both bad we
would have $h_2+h_3=h_1=h_2+h_4$, hence $h_3=h_4$, contradicting strictness. So one of them is
good.
\end{proof}

\begin{lemma}\label{lem:tricol}
Every triangle of $G(\mu)$ lies inside a single row or a single column of $\mu$.
\end{lemma}

\begin{proof}
By Theorem~\ref{thm:struct} each edge lies in a row or a column. Suppose a triangle $\{x,y,z\}$
has $x,y$ in row $i$ and $x,z$ in column $k$, with $x=(i,k)$, $y=(i,k_y)$, $z=(i_z,k)$,
$k_y\ne k$, $i_z\ne i$. Then $y$ and $z$ lie in different rows and different columns, so no
edge joins them. Hence all three edges are of the same kind and share a row, or share a column.
\end{proof}

\begin{lemma}\label{lem:row1}
Row~$1$ of $\mu$ contains no triangle if and only if $\mu_1\le2$, or $\mu_1=3$ and $n=2\ell+1$.
Dually, column~$1$ contains no triangle if and only if $\ell\le2$, or $\ell=3$ and
$n=2\mu_1+1$.
\end{lemma}

\begin{proof}
Apply Lemma~\ref{lem:line} with $r=\mu_1$ and $h_j=h_{1j}$. A triangle inside row~$1$ exists
iff a good pair exists. By Lemma~\ref{lem:line}(2) this fails iff $\mu_1\le2$, or $\mu_1=3$ and
$h_{12}+h_{13}=h_{11}$.

Suppose $\mu_1=3$. Then $h_{11}=\mu_1+\mu'_1-1=\ell+2$, $h_{12}=\mu_1-2+\mu'_2-1+1=\mu'_2+1$,
and $h_{13}=\mu_1-3+\mu'_3-1+1=\mu'_3$. Also $n=\mu'_1+\mu'_2+\mu'_3=\ell+\mu'_2+\mu'_3$
because $\mu_1=3$. Hence
\[
   h_{12}+h_{13}=h_{11}
   \iff \mu'_2+\mu'_3+1=\ell+2
   \iff (n-\ell)=\ell+1
   \iff n=2\ell+1 .
\]
The column statement is the conjugate.
\end{proof}

\begin{theorem}[classification of triangle-free $G(\mu)$]\label{thm:tri}
$G(\mu)$ contains no triangle if and only if $n\le3$ or $\mu\in\{(2,2),(3,2),(2,2,1)\}$.
\end{theorem}

\begin{proof}
$(\Leftarrow)$ For $n\le3$ and for the three listed $\mu$ every row and every column of $\mu$
has at most $3$ cells, and where it has exactly $3$ one checks $h_2+h_3=h_1$ directly; by
Lemmas~\ref{lem:line} and \ref{lem:tricol} there is no triangle. Explicitly: for
$\mu=(3)$ the row-$1$ hooks are $3,2,1$ with $2+1=3$; for $\mu=(2,2)$ and $\mu=(2,1)$ no line
has $3$ cells; for $\mu=(3,2)$ the row-$1$ hooks are $4,3,1$ with $3+1=4$ and every other line
has $\le2$ cells; $(1^3)$, $(1,1)$, $(2,2,1)$ are the conjugates of $(3)$, $(2)$, $(3,2)$.

$(\Rightarrow)$ Assume $n\ge4$ and $G(\mu)$ triangle-free. Then by Lemma~\ref{lem:row1} both
\[
  \bigl(\mu_1\le2 \ \text{ or }\ (\mu_1=3 \text{ and } n=2\ell+1)\bigr)
  \quad\text{and}\quad
  \bigl(\ell\le2 \ \text{ or }\ (\ell=3 \text{ and } n=2\mu_1+1)\bigr).
\]
In particular $\mu_1\le3$ and $\ell\le3$, so $n\le9$. Four cases.

\emph{(a) $\mu_1\le2$ and $\ell\le2$.} Then $n\le4$, so $n=4$ and $\mu=(2,2)$.

\emph{(b) $\mu_1\le2$, $\ell=3$, $n=2\mu_1+1$.} If $\mu_1=1$ then $n=3$, excluded. So
$\mu_1=2$, $n=5$, $\ell=3$: $\mu=(2,2,1)$.

\emph{(c) $\mu_1=3$, $n=2\ell+1$, $\ell\le2$.} If $\ell=1$ then $n=3$, excluded. So $\ell=2$,
$n=5$, $\mu_1=3$: $\mu=(3,2)$.

\emph{(d) $\mu_1=3$, $\ell=3$, $n=2\ell+1=7$.} Then $\mu\vdash7$ with $\mu_1=3$, $\ell=3$, so
$\mu=(3,3,1)$ or $\mu=(3,2,2)$. We rule both out using a \emph{second} line.
For $\mu=(3,3,1)$ we have $\mu'=(3,2,2)$ and the row-$2$ hooks are
\[
  h_{21}=3-1+3-2+1=4,\qquad h_{22}=3-2+2-2+1=2,\qquad h_{23}=3-3+2-2+1=1,
\]
so $h_{22}+h_{23}=3\notin\{4,2,1\}=H_2$ and $\{(2,1),(2,2),(2,3)\}$ is a triangle by
Lemma~\ref{lem:line}(1). The case $\mu=(3,2,2)=(3,3,1)'$ follows by conjugation (its
column-$2$ hooks are likewise $4,2,1$).

So the only triangle-free $\mu$ with $n\ge4$ are $(2,2)$, $(3,2)$, $(2,2,1)$.
\end{proof}

\begin{corollary}[(G2), in strengthened form]\label{cor:G2}
For every $\mu\vdash n$ with $n\ge4$ and $\mu\notin\{(2,2),(3,2),(2,2,1)\}$, the graph
$G(\mu)$ contains a triangle. Hence it contains an odd cycle, whose label product is $\ne1$ by
Lemma~\ref{lem:odd}. In particular this holds for every $\mu$ with $n\ge6$, which is (G2).
\end{corollary}

\begin{corollary}\label{cor:class}
\cite[Thm.~3.1]{Q140} is proved, for all $\mu$ and all $n$: the system \eqref{eq:local} is
consistent over $\Q(t)$ exactly for $n\le3$ \textup{(}all $\mu$\textup{)}, for $n=4$ with
$\mu=(2,2)$, and for $n=5$ with $\mu\in\{(3,2),(2,2,1)\}$; and inconsistent otherwise.
\end{corollary}

\begin{proof}
\emph{Inconsistency} for $n\ge4$, $\mu\notin\{(2,2),(3,2),(2,2,1)\}$: combine
Theorem~\ref{thm:G1}, Corollary~\ref{cor:G2} and Theorem~\ref{thm:graph} (whose hypothesis
Proposition~\ref{prop:reflect} is proved in \S\ref{sec:reflect}).

\emph{Consistency} in the nine remaining cases is proved by exhibiting a solution; each is
checked by substitution into \eqref{eq:local}, and each denominator below is a nonzero element
of $\Q(t)$. Listing $w$ against the columns $\gamma\in C(\mu)$ in the order shown:
\[
\begin{array}{ll}
\mu=(1): & w_{\varnothing}=1 \\
\mu=(2): & \tfrac{1}{t-1}\bigl(t,\,-1\bigr) \ \text{ on }\ \bigl((1),\varnothing\bigr) \\
\mu=(1^2): & \tfrac{1}{t-1}\bigl(-1,\,1\bigr) \ \text{ on }\ \bigl((1),\varnothing\bigr) \\
\mu=(3): & \tfrac{1}{2t-1}\bigl(t,\,t,\,-1\bigr) \ \text{ on }\ \bigl((2),(1),\varnothing\bigr) \\
\mu=(2,1): & \tfrac{1}{t^2-t+1}\bigl(1,\,t^2,\,-1\bigr) \ \text{ on }\ \bigl((2),(1,1),\varnothing\bigr) \\
\mu=(1^3): & \tfrac{1}{t(t-2)}\bigl(-t,\,-1,\,1\bigr) \ \text{ on }\ \bigl((1,1),(1),\varnothing\bigr) \\
\mu=(2,2): & \tfrac{1}{(t-1)^2}\bigl(-t,\,1,\,t,\,-1\bigr) \ \text{ on }\ \bigl((2,1),(2),(1,1),(1)\bigr) \\
\mu=(3,2): & \tfrac{1}{2t^2-2t+1}\bigl(-t,\,t^2,\,1,\,t,\,-1\bigr) \ \text{ on }\ \bigl((3,1),(2,2),(3),(1,1),(1)\bigr) \\
\mu=(2,2,1): & \tfrac{1}{t(t^2-2t+2)}\bigl(t,\,-t^2,\,t^2,\,1,\,-1\bigr) \ \text{ on }\ \bigl((2,2),(2,1,1),(1^3),(2),(1)\bigr)
\end{array}
\]
\vspace{-1.5em}
\end{proof}

\begin{remark}
Corollary~\ref{cor:class} does not rely on the $n\le9$ computation of \cite[\S6]{Q140}; that
computation now serves purely as a check. The hypothesis of \cite[Prop.~5.2]{Q140},
$\operatorname{rank}M^{(\mu)}(-1)=n$, is \emph{not} used anywhere above and remains open.
\end{remark}

\section{Correction to \cite[\S7(2)]{Q140}}\label{sec:correction}

That paper records, as an obstruction to finding triangles in row $1$:

\begin{quote}
``a type-(II) edge between two cells of row $i$ with hooks $d_1>d_2$ exists as soon as
$d_1+d_2>h_{i1}$ --- so the corner cell $(1,1)$ is joined to every other cell of row $1$. That
gives a hub but not, by itself, a triangle: for $\mu=(3,1,1,1,1)$ the row-$1$ hooks are $7,2,1$
and no second edge inside row $1$ exists.''
\end{quote}

The condition $d_1+d_2>h_{i1}$ is \emph{sufficient} but not necessary. By \eqref{eq:row} the
exact condition is $d_1+d_2\notin H_i$, and $d_1+d_2>h_{i1}=\max H_i$ is only one way for that
to happen.

\begin{example}[$\mu=(3,1,1,1,1)$]\label{ex:31111}
Here $n=7$, $\ell=5$, $\mu'=(5,1,1)$; the row-$1$ hooks are $H_1=\{7,2,1\}$. Then
\[
  h_{12}+h_{13}=2+1=3\notin\{7,2,1\}=H_1,
\]
so by \eqref{eq:row} the edge $(1,2)\sim(1,3)$ \textbf{exists}, and with the two hub edges of
Lemma~\ref{lem:hub} the set $\{(1,1),(1,2),(1,3)\}$ is a triangle. Its labels are
$-t^{5},-t^{5},-t$: on the abacus $b_1=2$, the holes below it are $-5<0<1$, and the three
relations read
\[
  w(2,0)=-t^{5}w(2,-5),\qquad w(2,1)=-t^{5}w(2,-5),\qquad w(2,1)=-t\,w(2,0).
\]
Substituting the first into the third and comparing with the second gives
$-t^{5}(1+t)\,w(2,-5)=0$, so $w(2,-5)=0$; with Theorem~\ref{thm:G1} the whole solution
vanishes. Consistently with Lemma~\ref{lem:odd}, the label product around the triangle is
$(-t^5)(-t)(-t^5)^{-1}=-t^{6}\ne1$.
\end{example}

This is not a repair of a false theorem: \cite[Thm.~6.3]{Q140} and \cite[Prop.~6.2]{Q140} are
correct, and the sentence quoted is a remark in the gaps section. But the remark was the stated
reason for not pursuing triangles, and it is void. Note also that the sufficient condition
suffices for the hub statement, so Lemma~\ref{lem:hub} --- and hence (G1) --- was already
within reach in \cite{Q140}; what was missing was the exact condition, which is what (G2) needs.

\section{Verification}\label{sec:verify}

All scripts are in \verb|proofs/scripts-2026-09-11-Q143/|.

\paragraph{Engines.} Three independent constructions of $G(\mu)$ were built and compared.
\begin{itemize}
\item \textbf{E1} (\emph{deriver}, \verb|graph.py|): from Proposition~\ref{prop:reflect}
      on the $\Z$-Maya abacus. Asserts $|V(G(\mu))|=n$ for every $\mu$ tested.
\item \textbf{E1b} (\emph{deriver}, \verb|hookgraph.py|): from Theorem~\ref{thm:struct},
      using hook lengths only --- no abacus.
\item \textbf{E2} (\emph{comparator / answer key}, \verb|matrix_engine.py|): builds
      $M^{(\mu)}(t)$ from rim-hook combinatorics on Young diagrams (connectivity and
      no-$2\times2$ tests on cell sets), then extracts the rows with exactly two nonzero
      entries. This uses neither the abacus nor hook lengths.
\end{itemize}
Results: \textbf{E1 $=$ E1b} on all $138$ partitions with $n\le10$ (edge sets identical).
\textbf{E1 $=$ E2} on all $66$ partitions with $n\le8$, agreeing on edge sets \emph{and} on
labels. This is sharper than the claim recorded in \cite{Q140}, which asserted only that the
two families span the same space.

\paragraph{Main table.} Per $n$, over all $\mu\vdash n$: number of $\mu$ with $G(\mu)$
connected; with $G(\mu)$ non-bipartite; for which the criterion of Theorem~\ref{thm:graph}
holds (tested on a spanning forest's fundamental cycles, which generate the cycle space).

\begin{center}
\begin{tabular}{rrrrrl}
\toprule
$n$ & $p(n)$ & connected & non-bipartite & criterion & triangle-free $\mu$ \\
\midrule
1 & 1  & 1  & 0  & 0  & $(1)$ \\
2 & 2  & 2  & 0  & 0  & $(2),(1^2)$ \\
3 & 3  & 3  & 0  & 0  & $(3),(2,1),(1^3)$ \\
4 & 5  & 5  & 4  & 4  & $(2,2)$ \\
5 & 7  & 7  & 5  & 5  & $(3,2),(2,2,1)$ \\
6 & 11 & 11 & 11 & 11 & --- \\
7 & 15 & 15 & 15 & 15 & --- \\
8 & 22 & 22 & 22 & 22 & --- \\
9 & 30 & 30 & 30 & 30 & --- \\
\midrule
$\Sigma$ & 96 & 96 & 87 & 87 & 9 \\
\bottomrule
\end{tabular}
\end{center}

The totals factorise as $96=1+2+3+5+7+11+15+22+30$ and $87=96-9$, the $9$ being the six
partitions with $n\le3$ together with $(2,2),(3,2),(2,2,1)$ --- exactly the exceptional set of
Theorem~\ref{thm:tri}, and exactly the set on which \eqref{eq:local} is consistent. (\cite{Q140}
reports ``$95$ partitions, $n\le9$''; that count omits $n=1$.) The columns
``non-bipartite'' and ``criterion'' agree entry by entry at every $n$: no even cycle with
nonzero signed exponent sum is ever needed.

\paragraph{The boundary population, explicitly.} A negative statement needs its
counterexamples exhibited. The nine triangle-free cases:
$(1),(2),(1^2),(3),(2,1),(1^3)$ give trees ($|E|=|V|-1$, so no cycle at all), while
\[
  (2,2):\ |V|=4,|E|=4;\qquad (3,2):\ |V|=5,|E|=5;\qquad (2,2,1):\ |V|=5,|E|=5,
\]
each having exactly one independent cycle, of \emph{even} length $4$ and signed exponent sum
$0$, hence label product $(-1)^4t^0=1$. For $(2,2)$ the cycle is
$(1,1)\to(2,1)\to(2,2)\to(1,2)\to(1,1)$ with exponents $0,+1,-0,-1$.

\paragraph{Checks of the proved statements.}
\begin{itemize}
\item Lemma~\ref{lem:hub}: all hub edges $(i,1)\sim(i,m)$ and $(1,k)\sim(i',k)$ present, for
      every $\mu$ with $n\le10$. $0$ failures.
\item Lemma~\ref{lem:line}: the predicate ``$r\ge4$, or $r=3$ and $h_2+h_3\ne h_1$'' agrees
      with direct search for a good pair on all $1624$ lines (rows and columns) of all
      partitions with $n\le11$. $0$ failures.
\item Lemma~\ref{lem:row1}: ``row~$1$ triangle-free $\iff\mu_1\le2$ or ($\mu_1=3$ and
      $n=2\ell+1$)'' holds for all $271$ partitions with $n\le12$. $0$ failures.
\item Theorem~\ref{thm:tri}: the classification agrees with a brute-force search over all
      $3$-subsets of cells (not using Lemma~\ref{lem:line}) on all $271$ partitions with
      $n\le12$. $0$ failures. Over $n\le10$, triangle-free $\iff$ bipartite.
\item Proposition~\ref{prop:reflect} (\S\ref{sec:reflect}): for every $\mu$ with $n\le8$
      and every configuration $b+c=b'+e$ of two beads and two holes, the row
      $\lambda$ of \eqref{eq:lam} was built and looked up in the matrix of engine E2. All
      $596$ configurations give $\lambda\ne\mu$ with $|\lambda|=n$, exactly two nonzero
      entries, in the predicted columns, with the predicted exponent difference. $0$
      failures. The count factorises as $596=298+298$ (type (I) / type (II), equal at every
      $n$ by conjugation) $=0+2+6+20+38+90+150+290$ over $n=1,\dots,8$, and equals the
      summed edge count $\sum_{\mu}|E(G(\mu))|$ --- each configuration contributes exactly
      one edge, with no coincidences.
\item Corollary~\ref{cor:class}: the nine solution vectors were obtained by solving
      $Mw=e_\mu$ over $\Q(t)$ with \textsc{SymPy} and then \emph{substituted back}; all nine
      satisfy $Mw-e_\mu=0$ identically. In the other direction, for every non-exceptional
      $\mu$ with $4\le n\le6$, $\operatorname{rank}[M\,|\,e_\mu]=\operatorname{rank}M+1$,
      confirming inconsistency independently of the graph criterion.
\item Example~\ref{ex:31111}: the edge $(1,2)\sim(1,3)$ of $G((3,1,1,1,1))$ is exhibited
      independently by E2 --- the row $\lambda=(2,2,1,1,1)$ of $M^{((3,1,1,1,1))}$ has exactly
      two nonzero entries, in columns $\gamma=(2,1,1,1,1)$ (entry $t^{0}$) and
      $\gamma=(1,1,1,1,1)$ (entry $t^{1}$), giving $w_{(1^5)}=-t^{-1}w_{(2,1^4)}$.
\end{itemize}

\paragraph{Negative controls.} Each perturbation must move a reported number; one that does
not would invalidate the control rather than confirm the claim.
\begin{enumerate}
\item \emph{Delete all type-(I) edges.} Connectivity collapses: over $n\le10$ the count of
      connected $G(\mu)$ falls from $96$ to $9$ (only $\mu=(n)$ survives at each $n$).
      \textbf{Moved.}
\item \emph{Delete all type-(II) edges.} Same totals, but by conjugation, not by being the
      same test: here only $\mu=(1^n)$ survives. Checked per $\mu$: control 1 on $\mu$ and
      control 2 on $\mu'$ agree in (\#components, non-bipartite) for all $\mu$, $n\le9$; and
      they differ on a single $\mu$, e.g.\ $\mu=(3,1)$ gives $2$ components versus $3$.
      \textbf{Moved.}
\item \emph{Replace the exact type-(II) condition \eqref{eq:row} by the sufficient condition
      $d_1+d_2>h_{i1}$ of \cite[\S7(2)]{Q140}.} The non-bipartite count moves at
      $n=4$ ($4\to3$), $n=5$ ($5\to4$) and $n=7$ ($15\to14$). \textbf{Moved} --- this is the
      control that isolates the error corrected in \S\ref{sec:correction}.
\item \emph{Drop the $r=3$ exception from Lemma~\ref{lem:line}}, i.e.\ predict ``a good pair
      exists iff $r\ge3$''. $56$ wrong predictions over the lines with $n\le11$, the smallest
      being row~$1$ of $(3)$ with hooks $3,2,1$; the correct form makes $0$. \textbf{Moved.}
\end{enumerate}

\section{What is \emph{not} proved}

\begin{enumerate}
\item The rank hypothesis $\operatorname{rank}M^{(\mu)}(-1)=n$ of \cite[Prop.~5.2]{Q140} ---
      linear independence of $\{p_{n-|\gamma|}s_\gamma:\gamma\in C(\mu)\}$ in $\Lam_n$ ---
      is untouched here. It is verified for $n\le8$ and is not used above.
\item Gap (2) of the earlier paper: Khanna--Loehr allow $T(\mu,L)$ to be an arbitrary subset
      of $R(n-L)$. Nothing here bears on that; Corollary~\ref{cor:class} refutes only the
      symmetric choice $T=S$.
\item \textbf{Completeness} of Proposition~\ref{prop:reflect}: the claim in
      \cite[\S6]{Q140} that (I) and (II) span the same space as \emph{all} two-term rows of
      $M$ is \emph{not} proved here --- \S\ref{sec:reflect} proves only that each (I)/(II)
      configuration is a genuine two-term row, which is all Theorem~\ref{thm:graph} needs.
      Completeness is verified in the stronger form ``the two sets of relations coincide,
      labels included'' for $n\le8$ (\S\ref{sec:verify}). It would be needed only for a
      converse to Theorem~\ref{thm:graph}.
\item Converse direction of Theorem~\ref{thm:graph}: failing the cycle criterion is not shown
      to imply consistency. The agreement of the last two columns of the table is a computed
      coincidence over $n\le9$, not a theorem. Corollary~\ref{cor:class} does not need it,
      since consistency in the nine exceptional cases is established by exhibiting solutions.
\end{enumerate}

\begin{thebibliography}{9}
\bibitem{Q140} Clio, \emph{A certificate family for the Khanna--Loehr local identity, and a
graph criterion for general $\mu$},
\verb|proofs/2026-09-10-c2-Q140-local-identity-certificate-family.tex|, 10 September 2026.
\end{thebibliography}

\end{document}
